\section{Background and Motivation}

In Viotti et al.~\cite{viotti2016consistency}, authors provided a
structured overview
of different consistency semantics appeared in distributed systems. Building
upon it's semantic definitions, we propose to use predicate logic derived from
those definitions to model consistency semantics, and reason about the
compositional properties of these semantics with SMT solvers (while Z3
  \cite{demoura2008z3} is used in our implementation, it's easily
  extensible with other
solver backends like CVC5 \cite{barbosa2022cvc5}).

To capture behaviors in distributed system executions, we model state
transitions on individual nodes using the notion of
\textit{operations}. Similar to
Viotti et al.~\cite{viotti2016consistency}, we define an operation as a tuple
containing: \texttt{proc} (process or node), \texttt{type} (read or
write), \texttt{obj} (the object
being acted upon), \texttt{ival} (input value), \texttt{oval} (output
value), \texttt{stime}
(start time), and \texttt{rtime} (return time). In our SMT encoding,
this tuple is
modeled as a sort, serving as the atomic building block for capturing system
behaviors and state transition constraints.

Partial and total orderings of operations are assigned through relation
functions defined in history and abstract
execution~\cite{burckhardt2014principles}
\cite{viotti2016consistency}. In the implementation, they are defined
as uninterpreted
functions with constraints~\cite{bryant2002modeling}. History
represent the set of all
operations executed in the system, while abstract execution instantiates/refines
the history by specifying which operations are visible to, or are arbitrated
over each other. By embedding these relations in first order logic, we can
leverage SMT solvers to check the satisfiability of a given configuration,
thereby determining whether a particular consistency model holds or if certain
axioms of that model are violated.

% TODO: summarize operations, axioms, and models
% NOTE: below shouldnt be a part of background

However, our current framework abstracts system-wide behavior from a single
node's perspective, applying constraints to operations primarily within the
scope of a single history and its corresponding abstract executions. While this
focuses the analysis and simplifies reasoning on local consistency properties
from the observing node's perspective, it also means that global behaviors
across multiple nodes need to be modeled separately on different set of symbols
to avoid conflicts in quantified variables. Further, to reason about cross
system behaviors, ``glue'' constraints need to be added to represent data flow
between different systems.
